\section{Implementation}

% TODO: write per Ch4 Rewrite Brief §4.1 — chapter intro paragraph(s) (V1/V2 framing)

\subsection{Implementation Overview}

The dashboard is deployed in two versions, both are fully featured and runnable in parallel from a single working tree. 
Version~1 (\texttt{code/}) is a Streamlit implementation that consists of two separate apps: one for instructors and one for lab assistants, each running their own Python process. 
Version~2 (\texttt{code2/}) is a single app, FastAPI process that servers both an asynchronous HTTP API at \texttt{/api/} and two buld React single-page applications (SPAs) for instructors and assistants respectively. 
V2 builds on V1, with every feature from V1 reimplemented in V2, additional features added on top and performance improvements. 
V1 is the first iteration of the dashboard, which served as a proof of concept and testbed for core features and design decisions.
V2 is the second iteration, which is a complete rewrite of V1 to address its limitations and scale up to a production-ready system. 


%A single analytical core sits behind both versions. The Python package \texttt{learning\_dashboard/} --- covering data ingestion, scoring, modelling, retrieval-augmented feedback, and saved-session persistence --- is reused by V2 unchanged, as a physical copy under \texttt{code2/learning\_dashboard/}. The only divergence between the two copies is a four-line refactor to the OpenAI-key guard that allows V2's FastAPI startup probe to verify dependencies without the application secrets present; no scoring code, no model code, and no retrieval code differs across the two versions. The two versions further share a single runtime-state file, \texttt{data/lab\_session.json}, coordinated through the \texttt{filelock} library: the live-session metadata, the joined-assistant roster, and the student-to-assistant assignment map are written and read by all three processes (V1 instructor, V1 assistant, V2 backend) under a five-second exclusive lock, so a session started in V1 becomes visible to V2 within approximately five seconds and vice versa. The implication is that the analytical contribution of this work is genuinely decoupled from the presentation layer, and the existence of V2 is itself evidence of that decoupling.

The motivation behind the second iteration of the dashboard is highlighted later on in the section. The first iteration had a Streamlit foundation that was ideal for rapid prototyping and early user feedback, but was limiting in terms of performance, scalability and flexibility. The limitation of rerunning the entire script on every interaction meant that the system struggled to maintain a responsive user experience which complicates asynchronous work. 
The second iteration addresses these limitations by moving these constraints to the presentation layer to React and the backend to FastAPI whilst keeping the analytical core untouched. 
The comparison between V1 and V2 is a recurring theme in the implementation section, as it provides a useful lens to explain the rationale behind design decisions and the trade-offs involved. 

This chapter is organised top-down so that Section \ref{sec:techstack} describes the tech stack used by both versions. 
Section \ref{sec:structure} describes the runtime processes, the shared runtime state and the deferred-actions pattern that underpins V1's reactive interactions. 
Section \ref{sec:pipeline} and \ref{sec:session} cover the data pipeline from the lab endpoint to the in-memory dataframe and the live and saved-session lifecycles. 
Section \ref{sec:analytics} covers the baseline analytics and Section \ref{sec:advanced} covers the advanced modelling layer of measurement confidence, item-response-theory difficulty, Bayesian knowledge-tracing mastery, and the improved struggle model that combines them. 
Section \ref{sec:rag} covers the retrieval-augmented feedback layer that produces real-time coaching suggestions for the lab assistant.
Section \ref{sec:views} and \ref{sec:assistant} walk through the instructor views and the lab-assistant flows respectively. 
Section \ref{sec:problems} highlights the implementation problems encountered and the resolutions adopted.


\subsection{Technology Stack} \label{sec:techstack}

% TODO: write per Ch4 Rewrite Brief §4.2 — drop in tab:techstack-v2 from Part H when ready

\subsection{System Structure} \label{sec:structure}

% TODO: write per Ch4 Rewrite Brief §4.3 — opening paragraph; reference fig:arch-v2 (architecture diagram)

\subsubsection{V1 - Instructor Process (Streamlit)} 

% TODO: write per Ch4 Rewrite Brief §4.3.1

\subsubsection{V1 - Assistant Process (Streamlit, mobile)} 
% TODO: write per Ch4 Rewrite Brief §4.3.2

\subsubsection{V2 - React + FastAPI} \label{sec:v2-react-fastapi}

% TODO: write per Ch4 Rewrite Brief §4.3.3 — open with the "Why V2?" paragraph naming concrete V1 limitations

\subsubsection{Shared Runtime State} 


% TODO: write per Ch4 Rewrite Brief §4.3.4

\subsubsection{Deferred-Actions Pattern}

% TODO: write per Ch4 Rewrite Brief §4.3.5

\subsection{Data Pipeline}

% TODO: write per Ch4 Rewrite Brief §4.4 — opening paragraph; reference fig:pipeline-flow

\subsubsection{Endpoint Retrieval and Parsing}

% TODO: write per Ch4 Rewrite Brief §4.4.1

\subsubsection{Data Normalisation and Structuring}

% TODO: write per Ch4 Rewrite Brief §4.4.2

\subsubsection{Caching Strategy}

% TODO: write per Ch4 Rewrite Brief §4.4.3

\subsection{Session Management}

% TODO: write per Ch4 Rewrite Brief §4.5

\subsubsection{Live Session Lifecycle}

% TODO: write per Ch4 Rewrite Brief §4.5.1

\subsubsection{Saved Session History and Restoration}

% TODO: write per Ch4 Rewrite Brief §4.5.2

\subsection{Analytics Implementation}

% TODO: write per Ch4 Rewrite Brief §4.6

\subsubsection{Incorrectness Scoring}

% TODO: write per Ch4 Rewrite Brief §4.6.1 — REQUIRED maths recap line (Marker-1 check)

\subsubsection{Baseline Student Struggle Model}

% TODO: write per Ch4 Rewrite Brief §4.6.2 — REQUIRED maths recap line; drop in tab:struggle-7sig

\subsubsection{Baseline Question Difficulty Model}

% TODO: write per Ch4 Rewrite Brief §4.6.3 — REQUIRED maths recap line; drop in tab:difficulty-5sig

\subsubsection{Collaborative Filtering}

% TODO: write per Ch4 Rewrite Brief §4.6.4 — REQUIRED maths recap line

\subsubsection{Mistake Clustering}

% TODO: write per Ch4 Rewrite Brief §4.6.5

\subsection{Advanced Model Implementation}

% TODO: write per Ch4 Rewrite Brief §4.7

\subsubsection{Measurement Confidence}

% TODO: write per Ch4 Rewrite Brief §4.7.1 — REQUIRED maths recap line

\subsubsection{Item Response Theory (IRT) Difficulty}

% TODO: write per Ch4 Rewrite Brief §4.7.2 — REQUIRED maths recap line; \cite{byrdLBFGSB1995}

\subsubsection{Bayesian Knowledge Tracing (BKT) Mastery}

% TODO: write per Ch4 Rewrite Brief §4.7.3 — REQUIRED maths recap line; \cite{corbettKnowledgeTracingModeling1995}; drop in tab:bkt-defaults

\subsubsection{Improved Struggle Model}

% TODO: write per Ch4 Rewrite Brief §4.7.4 — REQUIRED maths recap line

\subsection{RAG Suggested Feedback}

% TODO: write per Ch4 Rewrite Brief §4.8 — no inline supervisor attribution

\subsubsection{Hybrid Two-Layer Architecture}

% TODO: write per Ch4 Rewrite Brief §4.8.1 — include \cite{lewisRetrievalAugmentedGenerationKnowledgeIntensive2020}

\subsubsection{Embedding Model (all-MiniLM-L6-v2)}

% TODO: write per Ch4 Rewrite Brief §4.8.2

\subsubsection{Generation and Prompting}

% TODO: write per Ch4 Rewrite Brief §4.8.3

\subsubsection{Suggestion Caching}

% TODO: write per Ch4 Rewrite Brief §4.8.4

\subsubsection{Graceful Degradation and Privacy}

% TODO: write per Ch4 Rewrite Brief §4.8.5

\subsection{Instructor System Views}

% TODO: write per Ch4 Rewrite Brief §4.9.0 — opening paragraph + drop in tab:views-comparison (V1 vs V2 parity)

\subsubsection{In-Class View}

% TODO: write per Ch4 Rewrite Brief §4.9.1

\subsubsection{Student Detail View}

% TODO: write per Ch4 Rewrite Brief §4.9.2

\subsubsection{Question Detail View}

% TODO: write per Ch4 Rewrite Brief §4.9.3

\subsubsection{Data Analysis View}

% TODO: write per Ch4 Rewrite Brief §4.9.4

\subsubsection{Model Comparison View}

% TODO: write per Ch4 Rewrite Brief §4.9.5

\subsubsection{Settings View}

% TODO: write per Ch4 Rewrite Brief §4.9.6

\subsubsection{Previous Sessions View}

% TODO: write per Ch4 Rewrite Brief §4.9.7

\subsection{Lab Assistant System}

% TODO: write per Ch4 Rewrite Brief §4.10

\subsubsection{Session Join Flow}

% TODO: write per Ch4 Rewrite Brief §4.10.1

\subsubsection{Waiting and Assignment States}

% TODO: write per Ch4 Rewrite Brief §4.10.2

\subsubsection{Live Assistant Allocation and Mark-Helped}

% TODO: write per Ch4 Rewrite Brief §4.10.3

\subsection{Problems Encountered and Solutions}

% TODO: write per Ch4 Rewrite Brief §4.11 — drop in tab:problems
% Note: sound-effects discussion now lives inside §4.9.6 Settings View (the toggle is there); old standalone §4.11 Sound Effects subsection removed during 2026-05-20 layout refinement.

% Old V1-prototype draft prose has been moved to implementation-old.tex (sidecar, not \include'd in main.tex). Decision 2026-05-20: nothing in there is being reused.
